\documentclass[11pt,a4paper]{article}
\usepackage{fontspec,amsmath,amssymb,amsthm,microtype}
\setmainfont{lmroman10-regular.otf}[BoldFont=lmroman10-bold.otf,ItalicFont=lmroman10-italic.otf,BoldItalicFont=lmroman10-bolditalic.otf]
\setmonofont{lmmono10-regular.otf}[Scale=MatchLowercase]
\usepackage[margin=25mm]{geometry}
\usepackage[unicode,hidelinks]{hyperref}
\setlength{\emergencystretch}{3em}
\setlength{\parskip}{0.4em}
\setlength{\parindent}{0pt}
\allowdisplaybreaks[2]
\raggedbottom
\title{\texorpdfstring{One-\allowbreak{}dimensional Prékopa–Leindler by constructed monotone transport}{One-dimensional Prékopa–Leindler by constructed monotone transport}}
\author{LangevinLean mathematical exposition}
\date{}
\begin{document}\maketitle
\noindent Source module: \texttt{IPLA/\allowbreak{}Prekopa\allowbreak{}Leindler.\allowbreak{}lean}. The accompanying HTML pairs this mathematical exposition with the exact Lean statements and records the scope of the formalization.


The proof constructs transport directly from cumulative integrals of positive continuous densities. It derives the one-dimensional Prékopa–Leindler inequality without postulating a transport map or marginal log-concavity theorem.


\section{\texorpdfstring{Cumulative integrals are differentiable and strictly increasing}{Cumulative integrals are differentiable and strictly increasing}}


A density here is a continuous strictly positive real function \(f\), integrable on \(\mathbb R\), with \(\int f=1\). Define \(F(x)=\int_{(-\infty,x]}f\). The difference formula

\[F(y)-F(0)=\int_0^y f(t)dt\]

and the fundamental theorem of calculus give the strict derivative \(F'(x)=f(x)\). Differentiability implies continuity, and positivity of the derivative implies strict monotonicity. The density assumptions are precisely the inputs required for these assertions.




\textbf{Correspondence with the displayed declarations.}


\begin{itemize}

\item \texttt{Density}: A positive, continuous probability density on the real line.
\item \texttt{cumulative}: Cumulative integral of an ordinary real density.
\end{itemize}


{\raggedright\noindent Lean declarations: \texttt{Density}, \texttt{cumulative}, \texttt{cumulative\_\allowbreak{}has\allowbreak{}Strict\allowbreak{}Deriv\allowbreak{}At}, \texttt{cumulative\_\allowbreak{}continuous}, \texttt{cumulative\_\allowbreak{}strict\allowbreak{}Mono}.\par}

\section{\texorpdfstring{The cumulative range is exactly the open unit interval}{The cumulative range is exactly the open unit interval}}


Integrability implies \(F(x)\to0\) as \(x\to-\infty\). Splitting the total integral at \(x\) gives \(F(x)=1-\int_{(x,\infty)}f\), so \(F(x)\to1\) as \(x\to+\infty\). Nonnegativity gives \(0\le F\le1\); strict monotonicity, comparing \(x\) with \(x-1\) and \(x+1\), improves this to \(0<F(x)<1\). For every \(u\in(0,1)\), the endpoint limits supply points on either side of \(u\), and the intermediate value theorem supplies \(F(x)=u\). Therefore \(F(\mathbb R)=(0,1)\).



{\raggedright\noindent Lean declarations: \texttt{cumulative\_\allowbreak{}tendsto\_\allowbreak{}at\allowbreak{}Bot}, \texttt{cumulative\_\allowbreak{}tendsto\_\allowbreak{}at\allowbreak{}Top}, \texttt{cumulative\_\allowbreak{}mem\_\allowbreak{}Ioo}, \texttt{cumulative\_\allowbreak{}range}.\par}

\section{\texorpdfstring{Construct and differentiate the monotone transport}{Construct and differentiate the monotone transport}}


For two densities with cumulative maps \(F,G\), set \(T=G^{-1}\circ F\), using the inverse function on its range. Since both ranges are \((0,1)\), \(G(T(x))=F(x)\) everywhere. Strict monotonicity makes \(G\) injective. The one-dimensional inverse function theorem applies because \(G'(T(x))=g(T(x))>0\), and the chain rule yields

\[T'(x)=\frac{f(x)}{g(T(x))}>0.\]

Thus \(T\) is strictly increasing, and the density ratio is a proved derivative of an explicitly constructed function.




For positive continuous unit-mass densities, \(F(x)=\int_{-\infty}^xf\) and \(G(y)=\int_{-\infty}^yg\) are differentiable with strictly positive derivatives, hence strictly increasing, and their limits give range \((0,1)\). Define \(T=G^{-1}\circ F\), so \(G(T(x))=F(x)\). The inverse-function derivative gives \(T'(x)=f(x)/g(T(x))>0\). The cumulative identities and strict monotonicity declarations justify this change of variables on the whole real line.




\textbf{Correspondence with the displayed declarations.}


\begin{itemize}

\item \texttt{transport}: The monotone transport from density \texttt{f} to density \texttt{g}.
\end{itemize}


{\raggedright\noindent Lean declarations: \texttt{transport}, \texttt{cumulative\_\allowbreak{}transport}, \texttt{transport\_\allowbreak{}has\allowbreak{}Strict\allowbreak{}Deriv\allowbreak{}At}, \texttt{transport\_\allowbreak{}strict\allowbreak{}Mono}.\par}

\section{\texorpdfstring{The normalized inequality is a change-\allowbreak{}of-\allowbreak{}variables calculation}{The normalized inequality is a change-of-variables calculation}}


Assume \(0\le t\le1\) and \(h((1-t)x+ty)\ge f(x)^{1-t}g(y)^t\). Put \(S(x)=(1-t)x+tT(x)\) and \(r=f(x)/g(T(x))>0\). Then \(S'=1-t+tr>0\), and convexity of the exponential gives \(r^t\le1-t+tr=S'\). At \(y=T(x)\), the pointwise premise therefore implies

\[f(x)=f(x)^{1-t}g(T(x))^t r^t\le h(S(x))S'(x).\]

Since \(S\) is strictly increasing, the one-dimensional Jacobian theorem changes the right-hand integral to the integral of \(h\) over \(S(\mathbb R)\). Nonnegativity then bounds it by the integral over all \(\mathbb R\), proving \(1\le\int h\). The formal conclusion uses an extended nonnegative integral, so \(h\) need not itself be integrable or continuous.




Set \(S(x)=(1-t)x+tT(x)\). The interpolation hypothesis bounds \(h(S(x))\) below by \(f(x)^{1-t}g(T(x))^t\). Weighted arithmetic–geometric mean gives \(S'(x)=(1-t)+tT'(x)\ge T'(x)^t\). Since \(T'=f/g(T)\), multiplication yields \(h(S(x))S'(x)\ge f(x)\). The strictly increasing change-of-variables formula then proves \(\int h\ge\int f=1\). For arbitrary positive finite masses, divide \(f,g\) by their integrals, apply the normalized statement, and multiply back; logarithmic mass terms recover \((\int f)^{1-t}(\int g)^t\).




\textbf{Correspondence with the displayed declarations.}


\begin{itemize}

\item \texttt{normalized\_\allowbreak{}prekopa\_\allowbreak{}leindler}: The normalized one-dimensional Prékopa--Leindler inequality, proved by the constructed monotone transport. The third weight need not be continuous or integrable; the conclusion is an extended nonnegative integral.
\end{itemize}


{\raggedright\noindent Lean declarations: \texttt{normalized\_\allowbreak{}prekopa\_\allowbreak{}leindler}.\par}

\section{\texorpdfstring{Remove the unit-\allowbreak{}integral normalization}{Remove the unit-integral normalization}}


An everywhere positive integrable function on a nonzero measure space has strictly positive integral, since its support has positive measure. For the Prékopa weights write \(A=\int f>0\), \(B=\int g>0\), and \(C=A^{1-t}B^t>0\). Apply the normalized inequality to \(f/A\), \(g/B\), and \(h/C\). Expanding \(\log(f/A)=\log f-\log A\) verifies its pointwise premise. If \(h\) is nonnegative and integrable, the extended conclusion converts to \(1\le(\int h)/C\). Multiplication by \(C\) gives

\[\exp\!\left((1-t)\log\int f+t\log\int g\right)\le\int h.\]

The positivity lemma also applies to the zero-dimensional product measure needed later.




\textbf{Correspondence with the displayed declarations.}


\begin{itemize}

\item \texttt{integral\_\allowbreak{}positive}: Strict positivity of the real integral of an everywhere positive integrable function, including the zero-dimensional product space used in tensorization.
\item \texttt{prekopa\_\allowbreak{}leindler\_\allowbreak{}real}: Prékopa--Leindler on the real line for arbitrary positive continuous integrable first and second weights. No normalization is assumed.
\end{itemize}


{\raggedright\noindent Lean declarations: \texttt{integral\_\allowbreak{}positive}, \texttt{prekopa\_\allowbreak{}leindler\_\allowbreak{}real}.\par}

\end{document}
